\documentclass[11pt]{article}
\usepackage{preamble}
\setlength{\parskip}{4pt}

\begin{document}

\daybanner{AP Statistics \textperiodcentered\ Unit 2 \textperiodcentered\ Topic 2.8 \textperiodcentered\ B: 2026-10-30 \textperiodcentered\ E: 2026-11-20 \textperiodcentered\ TEACHER KEY}{Average Payoff, Typical Ticket}

\sectionbanner{CED TETHER}

{\small
\begin{itemize}[leftmargin=1.5em,itemsep=2pt,topsep=2pt]
  \item \textbf{Skill 2.B} --- Construct numerical or graphical representations of distributions.
  \item \textbf{Skill 4.B} --- Interpret statistical calculations and findings to assign meaning or assess a claim.
  \item \textbf{EU VAR-5} --- Probability distributions may be used to model variation in populations.
  \item \textbf{LO VAR-5.A} --- Represent the probability distribution for a discrete random variable. [Skill 2.B]
  \item \textbf{EK VAR-5.A.1} --- The values of a random variable are the numerical outcomes of random behavior.
  \item \textbf{EK VAR-5.A.2} --- A discrete random variable is a variable that can only take a countable number of values. Each value has a probability associated with it. The sum of the probabilities over all of the possible values must be 1.
  \item \textbf{EK VAR-5.A.3} --- A probability distribution can be represented as a graph, table, or function showing the probabilities associated with values of a random variable.
  \item \textbf{EK VAR-5.A.4} --- A cumulative probability distribution can be represented as a table or function showing the probability of being less than or equal to each value of the random variable.
  \item \textbf{LO VAR-5.B} --- Interpret a probability distribution. [Skill 4.B]
  \item \textbf{EK VAR-5.B.1} --- An interpretation of a probability distribution provides information about the shape, center, and spread of a population and allows one to make conclusions about the population of interest.
\end{itemize}
}
\textbf{Essential question:} How can an average payout and the most common outcome describe different parts of the same distribution?\par
\textbf{DOK-3 skill:} 4.B \textperiodcentered\ \textbf{FRQ pattern:} {\small\texttt{expected-\allowbreak{}value-\allowbreak{}vs-\allowbreak{}typical-\allowbreak{}outcome-\allowbreak{}what-\allowbreak{}the-\allowbreak{}mean-\allowbreak{}tells-\allowbreak{}one-\allowbreak{}person}}\par
\textbf{DOK rationale:} Part (c) weighs competing claims using the day's numerical or graphical evidence and explains a limitation or consequence; the judgment requires linked reasoning beyond a calculation.\par

\frameworkphaseheader{Do Now (first take) $\rightarrow$ Explore (video + follow-along) $\rightarrow$ Exit (finish + turn in)}{1 $\rightarrow$ 3}{45}{%
\begin{itemize}[leftmargin=1.5em,itemsep=2pt,topsep=2pt]
  \item Show the board at the bell and collect a one-sentence first take without confirming a claim.
  \item Run the named video follow-along, then allow about 10 minutes for parts (a)--(c).
  \item Ask for evidence linked to a claim. Collect the sheet and hand-score part (c) with E/P/I.
\end{itemize}
}{%
\begin{itemize}[leftmargin=1.5em,itemsep=2pt,topsep=2pt]
  \item Compare the given average with the ticket distribution, finish (a)--(c), and turn in the sheet.
\end{itemize}
}{%
\begin{itemize}[leftmargin=1.5em,itemsep=2pt,topsep=2pt]
  \item Which number or visible feature supports your claim?
  \item Why does the other claim go beyond that evidence?
  \item What would you tell someone who chose the other claim?
\end{itemize}
}{Circulate and ask students to connect their choice to evidence. Score part (c) using the three required reasoning elements; do not require dropped extension work.}

\begin{callout}{\IconWarn\ WATCH FOR}{calloutred}
\begin{itemize}[leftmargin=1.5em,itemsep=2pt,topsep=2pt]
  \item A choice with copied numbers but no explanation of how they support it.
  \item Treating a model or average as a guaranteed individual outcome.
\end{itemize}
\end{callout}

\sectionbanner{SCORING PART (c) --- E / P / I}

\textbf{Expected elements}
\begin{itemize}[leftmargin=1.5em,itemsep=2pt,topsep=2pt]
  \item \textbf{reasoning-1} (required) --- Reconciles the given average payout with most tickets paying nothing
  \item \textbf{reasoning-2} (required) --- Uses the 2-dollar cost and 87.5-percent loss chance to assess Priya's claim
  \item \textbf{reasoning-3} (required) --- Distinguishes an average from the possible result of one ticket
\end{itemize}

{\small\renewcommand{\arraystretch}{1.25}
\noindent\begin{tabular}{|p{0.5in}|p{5.9in}|}
\hline
\textbf{E} & Includes all three required reasoning elements above, with correct contextual evidence. \\ \hline
\textbf{P} & Includes at least one substantive correct reasoning link above, but omits or misstates another required element. \\ \hline
\textbf{I} & Includes no substantive correct reasoning link above; an unsupported choice or copied number alone is insufficient. \\ \hline
\end{tabular}}

\textbf{Common mistakes}
\begin{itemize}[leftmargin=1.5em,itemsep=2pt,topsep=2pt]
  \item Treating the average payout as the most common ticket outcome
  \item Ignoring the 2-dollar ticket cost
  \item Saying most tickets paying nothing makes a positive average payout impossible
\end{itemize}

\clearpage
\focusbanner{TODAY'S PROBLEM --- annotated}

Suppose a club sells 200 raffle tickets for \$2 each. One ticket wins \$100, four tickets win \$25 each, 20 tickets win \$5 each, and the other tickets win nothing. Let $X$ be the buyer's net profit from one randomly chosen ticket: prize minus ticket cost.\par\smallskip \textbf{Priya:} \emph{The average gross payout is \$1.50 per ticket, so buying one is a good deal.}\par \textbf{Mateo:} \emph{Most tickets win nothing.}\par
\begin{center}
\textbf{Ticket net-profit distribution}\par\smallskip
\begin{tabular}{|l|c|c|}
\hline
\textbf{Prize} & \textbf{Net $x$} & \textbf{$P(X=x)$} \\ \hline
\textbf{Top prize} & $98$ & $1/200$ \\ \hline
\textbf{Second prize} & $23$ & $4/200$ \\ \hline
\textbf{Small prize} & $3$ & $20/200$ \\ \hline
\textbf{No prize} & $-2$ & $175/200$ \\ \hline
\end{tabular}
\end{center}
\begin{firsttakebox}{FIRST TAKE (5 min)}
Would you buy one ticket if your goal were to make money? Use one number or outcome from the problem.\par
\answer{Not graded. Students may focus on the top prize or on the 175 losing tickets; the distribution should connect those views.}
\end{firsttakebox}

\textit{Rules callout ``DISTRIBUTIONS AND ONE OUTCOME (keep on the page)'' is printed on the student sheet.}\par\medskip

\dokbadge{1}~\textbf{(a)} State the four possible net profits $X$. Which is most likely?\par
\answer{The possible net profits are $98$, $23$, $3$, and $-2$ dollars. The most likely value is $X=-2$, because 175 of the 200 tickets pay no prize.}

\dokbadge{2}~\textbf{(b)} Check that the listed probabilities add to 1. Find the chance that a ticket loses money.\par
\answer{$(1+4+20+175)/200=1$, with every probability between 0 and 1. Only the no-prize ticket loses money, so $P(X<0)=175/200=0.875=87.5\%$.}

\dokbadge{3}~\textbf{(c)} Can Priya's average payout and Mateo's statement both be true? Use the distribution to explain. Decide whether Priya's good-deal claim follows for a buyer trying to make money, using the ticket cost and possible net profits. Write 3--4 sentences.\par
\answer{Both statements can be true: a few large prizes raise the average payout even though 175 of 200 tickets pay nothing. An average does not promise that one buyer receives \$1.50; the only net profits are \$98, \$23, \$3, and $-\$2$. Priya's claim does not follow for making money because the \$1.50 average payout is below the \$2 cost and the chance of losing \$2 is $87.5\%$. A buyer could still choose to support the club, but that is a different reason from expecting to make money.}

\begin{summaryexitbox}{EXIT}
One thing the video changed about my first take: \hrulefill
\end{summaryexitbox}

\end{document}
